\begin{algorithm}[H]
\caption{\ExpectedTree: the polynomial expected-depth algorithm}
\label{alg:expected-tree}
\KwIn{State $S$, original positive rational priors;
calibrated $\beta,\delta,p_*$.}
\KwOut{A supplied-query identification tree.}
\If{$|S|\le1$}{\KwRet its leaf\;}
\If{a mass-balanced supplied line $\ell$ exists}{
 \KwRet $\ell$ with children $\ExpectedTree(S_\ell^-)$
 and $\ExpectedTree(S_\ell^+)$\;
}
Compute $\HullData(S)$, its pattern classes, and zero weight $p_0=w(A_0)$\;
\eIf{there are at most two real extreme hull tests}{
 Build their direct full-pattern stage\;
}{
 Use edges of the full hull as blocks, retaining its free origin marker\;
 Select the schedule and stored modes by Algorithm~\ref{alg:two-phase-family}\;
 Build discovery along its all-negative prefix, with those modes on each first hit\;
 At each positive exit build expected $\EndpointCode$ to recover the full pattern\;
}
Attach $\ExpectedTree(A)$ to each nonzero pattern class $A$\;
Use a zero-query leaf for the zero singleton; omit it when absent\;
\KwRet the tree after contracting constant and cached questions\;
\end{algorithm}

The cases in the algorithm have a direct purpose: singleton states
already identify their hypothesis, balanced lines split immediately,
and flat states use the hull schedule to discover a full pattern before
recursing within its class. Each recursive call reuses the fixed
calibration and full supplied-query family. Since the incidence table
and prior weights are known, schedule scores, endpoint modes, and child
codes can all be determined before the hidden hypothesis is queried;
the routine therefore returns a concrete tree for it to follow.
